% =============================================================================
% SECCIÓN 2: MATERIALES Y MÉTODOS
% Capítulo 2: Instrumentación IoT y Caracterización Metrológica
% =============================================================================


% -----------------------------------------------------------------------------
% 2.1. DISEÑO EXPERIMENTAL Y SUSTRATOS
% -----------------------------------------------------------------------------
\subsection{Diseño experimental y caracterización de sustratos}
\label{subsec:cap2_diseno_sustratos}

El diseño experimental se estructuró como un ensayo de laboratorio
controlado con arreglo de bloques completos al azar, evaluando dos
tipos de sustratos agroindustriales del Valle del Cauca, Colombia:
la dieta D1 (alta carga proteica, base pollo) y la dieta D4 (alta
carga fibrosa, base vegetal). Cada tratamiento se replicó en tres
corridas independientes (réplicas biológicas), con un total de seis
ciclos experimentales de bioconversión de 15 días cada uno.

Las larvas neonatas de \textit{H. illucens} (estadio L1, 3--5 días
post-eclosión) se distribuyeron en repositorios internos de
polipropileno alimenticio (20 $\times$ 15 $\times$ 10 cm) con una
densidad de siembra de 0.5 g larva$\cdot$cm$^{-2}$, siguiendo los
rangos óptimos reportados para maximizar la tasa de asimilación
\autocite{nayakHermetiaIllucensDiptera2023}. El sustrato se
acondicionó a una humedad inicial del 65--70\,\% (base húmeda) y se
mantuvo en cámara de crecimiento a $30 \pm 2$\,°C y HR del 60--70\,\%,
condiciones que favorecen el desarrollo metabólico máximo según
\autocite{salamEffectDifferentEnvironmental2022}.

La caracterización proximal de los sustratos incluyó: contenido de
humedad (deshidratación a 105\,°C hasta peso constante), cenizas
(incineración a 550\,°C), proteína bruta (método Kjeldahl, factor
6.25), lípidos totales (extracción Soxhlet con hexano), fibra bruta
(digestión ácido-base secuencial) y carbono/nitrógeno total (analizador
elemental CHNS-O). El pH se midió en suspensión 1:5 (sustrato:agua
destilada) con electrodo de vidrio calibrado.

% -----------------------------------------------------------------------------
% 2.2. ARQUITECTURA DEL NODO IoT
% -----------------------------------------------------------------------------
\subsection{Arquitectura del sistema de adquisición}
\label{subsec:cap2_arquitectura_iot}

El sistema de monitoreo implementó una arquitectura de tres capas:
percepción, procesamiento edge y telemetría, adaptada de los esquemas
modulares propuestos por \autocite{vanIntegrationInternetofThingsSustainable2022}
y \autocite{duobieneDevelopmentWirelessSensor2022}.

\textbf{Capa de percepción.} El nodo sensor se construyó sobre un
microcontrolador ESP32-WROOM-32 (Espressif Systems), seleccionado por
su capacidad de procesamiento dual-core (240\,MHz), conectividad WiFi
integrada y bajo consumo energético en modo deep-sleep (10\,$\mu$A).
El bus de sensores incluyó:

\begin{itemize}
    \item \textbf{CO$_2$ y CH$_4$:} Sensores NDIR de bajo costo
    (modelos SCD41 para CO$_2$, Sensirion; y un sensor NDIR dual
    calibrado de fábrica para CH$_4$ en rango 0--5000\,ppm). Ambos
    operan con resolución de 1\,ppm, precisión de $\pm(30 + 3\%)$\,ppm
    y deriva máxima anual de $\pm10$\,ppm para CO$_2$.
    
    \item \textbf{Variables ambientales:} Sensor SHT3x (Sensirion)
    para temperatura ($-40$ a 125\,°C, $\pm0.3$\,°C) y humedad
    relativa (0--100\,\%, $\pm2$\,\%).
    
    \item \textbf{Condicionamiento:} Cápsula de PTFE hidrofóbica
    (membrana ePTFE, poro 0.2\,$\mu$m) sobre la óptica NDIR para
    prevenir condensación directa sin inducir retardo de difusión
    significativo ($\tau_{63} < 30$\,s).
\end{itemize}

\textbf{Capa de procesamiento edge.} El ESP32 ejecutó un firmware
desarrollado en C++ (framework Arduino) con las siguientes rutinas:
(i) lectura multiplexada de sensores cada 30\,s; (ii) filtro digital
de media móvil exponencial (EMA, $\alpha = 0.3$) para atenuación de
ruido de alta frecuencia; (iii) almacenamiento temporal en buffer
circular de memoria Flash (4\,MB) ante fallos de conectividad; e
(iv) transmisión asincrónica vía MQTT (protocolo ligero) hacia un
broker EMQX desplegado en instancia local.

\textbf{Capa de contención experimental.} Para desacoplar la unidad
biológica de la unidad de medición, se empleó un sistema de doble
contención: el repositorio de cría (unidad biológica) se insertó en
una cámara de medición hermética de acrílico transparente (volumen
interno 5.2\,L) equipada con los sensores NDIR en el espacio de
cabeza. El protocolo de medición siguió un ciclo de
\textit{Cierre--Ventilación--Reclosure}:
(1) inserción del repositorio y cierre hermético, registrando
concentraciones hasta saturación o estabilización;
(2) apertura de la cámara y extracción del repositorio para
re-oxigenación pasiva del sustrato (10\,min); y
(3) reintroducción y segundo cierre hermético para evaluar la
histeresis metabólica.

% -----------------------------------------------------------------------------
% 2.3. CALIBRACIÓN Y COMPENSACIÓN CRUZADA
% -----------------------------------------------------------------------------
\subsection{Protocolo de calibración y compensación cruzada}
\label{subsec:cap2_calibracion}

\textbf{Calibración estática.} Previo al despliegue experimental, los
sensores NDIR se sometieron a calibración de dos puntos con gases
patrón certificados (CO$_2$: 0\,ppm aire sintético cero; 5000\,ppm
mezcla N$_2$/CO$_2$, incertidumbre $\pm2$\,\%, trazable a NIST).
La calibración se ejecutó dentro de la cámara de medición a
25\,°C y 50\,\% HR para replicar las condiciones de referencia del
fabricante.

\textbf{Compensación cruzada por temperatura y humedad.} Dado que
las lecturas NDIR exhiben deriva térmica y sensibilidad a la
humedad relativa \autocite{biagiDevelopmentMachineLearningbased2024},
se implementó un modelo de corrección multivariado en el firmware
del nodo:

\begin{equation}
    C_{\text{corr}} = C_{\text{raw}} \cdot \frac{T_{\text{cal}}}{T_{\text{amb}}} \cdot \frac{P_{\text{amb}}}{P_{\text{cal}}} \cdot \left[ 1 + \beta \cdot (\text{RH}_{\text{amb}} - \text{RH}_{\text{cal}}) \right] + \gamma \cdot (T_{\text{amb}} - T_{\text{cal}})
    \label{eq:compensacion_ndir}
\end{equation}

donde $C_{\text{raw}}$ es la concentración bruta (ppm),
$T_{\text{cal}} = 298.15$\,K, $P_{\text{cal}} = 101.325$\,kPa,
$\text{RH}_{\text{cal}} = 50$\,\%, $\beta = 2.1 \times 10^{-4}$
(\%RH)$^{-1}$ (coeficiente de sensibilidad a la humedad estimado
empíricamente para el rango 60--90\,\% HR), y $\gamma = -0.45$\,ppm\,K$^{-1}$
(coeficiente térmico de deriva).

\textbf{Validación dinámica.} Durante las corridas experimentales,
se tomaron muestras de aire del espacio de cabeza (10\,mL) mediante
jeringa hermética en tres instantes por ciclo (t = 0, 7 y 14 días),
analizadas por cromatografía de gases (GC-TCD, Agilent 7890B)
como método de referencia. Los pares $(C_{\text{sensor}},
C_{\text{GC}})$ se usaron para ajuste recursivo de los coeficientes
$\beta$ y $\gamma$ mediante regresión lineal multivariada.

% -----------------------------------------------------------------------------
% 2.4. MÉTRICAS DE VALIDACIÓN
% -----------------------------------------------------------------------------
\subsection{Métricas de desempeño metrológico}
\label{subsec:cap2_metricas}

El desempeño del sistema se evaluó mediante cuatro indicadores
cuantitativos, alineados con la hipótesis H1:

\begin{enumerate}
    \item \textbf{Coeficiente de correlación de Pearson} ($R$):
    linealidad entre la señal del sensor y el método de referencia.
    Umbral de aceptación: $R > 0.85$.
    
    \item \textbf{Error Porcentual Absoluto Medio} (MAPE):
    \begin{equation}
\text{MAPE} = \frac{100}{N} \sum_{i=1}^{N} \left| \frac{C_{\text{sensor},i} - C_{\text{ref},i}}{C_{\text{ref},i}} \right| \quad (\%)
        \label{eq:mape}
    \end{equation}
    Umbral de aceptación: MAPE $\leq 15$\,\%.
    
    \item \textbf{Raíz del Error Cuadrático Medio} (RMSE):
    \begin{equation}
        \text{RMSE} = \sqrt{\frac{1}{N} \sum_{i=1}^{N} (C_{\text{sensor},i} - C_{\text{ref},i})^2}
        \label{eq:rmse}
    \end{equation}
    
    \item \textbf{Incertidumbre expandida} ($U$, $k=2$):
    calculada conforme a la propagación de errores tipo A (repetibilidad)
    y tipo B (resolución del sensor, deriva del patrón, coeficientes de
    compensación), combinada como:
    \begin{equation}
        u_c = \sqrt{u_A^2 + u_B^2}
        \label{eq:incertidumbre_combinada}
    \end{equation}
    y reportada como $U = 2 \cdot u_c$ para un nivel de confianza del
    95\,\%.
\end{enumerate}

Adicionalmente, se reportó la \textbf{tasa de pérdida de paquetes}
de datos (packet loss) durante la telemetría MQTT y la
\textbf{continuidad operativa} del nodo (horas de registro efectivo
vs. horas programadas) como indicadores de robustez del sistema IoT.